% GENERATED FILE. Edit canonical lessons, then run npm run generate:books.
% canonical-source-hash: fnv1a64:74b8028b9c2e8693
% canonical-lessons: PT-C328-navegador, PT-C328-caravela, PT-C328-naufragar, PT-C328-naufragio, PT-C328-tesouro

\chapter{Pelo mar fora}
\label{ch:pt-pelo-mar-fora}
\hlchaptermodality{328}

\begin{chapteropening}
\textbf{By the end of this chapter:} \emph{I can tell a sea story: sailing, navigators, caravels, shipwrecks and treasure.}

Complete the last lesson of chapter 328: I can tell a sea story: sailing, navigators, caravels, shipwrecks and treasure.
\end{chapteropening}

\section[o navegador]{o navegador — the navigator, the seafarer; the web browser}

\label{lesson:PT-C328-navegador}



\begin{quote}
Before the new one: say the Portuguese for the throne, then the Portuguese for the crown.
\end{quote}

\subsection*{You'll want to know: o navegador}
\textbf{o navegador} — \textquotedblleft{}the navigator, the seafarer; the web browser\textquotedblright{}.

\begin{quote}
O navegador desenhou um mapa da costa.
\end{quote}

\begin{quote}
\emph{The navigator drew a map of the coast.}
\end{quote}

On a computer, \emph{o navegador} is the web browser.

\begin{cousinweb}
\emph{navegador} comes from Latin \emph{nāvigātor} \textquotedblleft{}a sailor\textquotedblright{} (EtymDB), from \emph{nāvigāre} \textquotedblleft{}to sail\textquotedblright{}.
\end{cousinweb}

\begin{grammarlens}[title={What you've built}]
One more word for telling, explaining and describing at B1.
\end{grammarlens}

\subsection*{Your turn}
\begin{itemize}
\raggedright
  \item \emph{Say it:} \emph{o navegador}
  \item \emph{Say it:} \emph{o navegador}, once more
  \item \emph{Say it:} a sentence of your own with \emph{o navegador}
  \item \emph{Recall:} say the Portuguese for the throne, then the Portuguese for the crown, then say \emph{o navegador} again
\end{itemize}

\begin{tcolorbox}[breakable,colback=teal!4,colframe=teal!35!black,title={Before you move on}]
What does o navegador mean? (\textquotedblleft{}The navigator, the seafarer; the web browser\textquotedblright{}.) Say it once more.
\end{tcolorbox}
\section[a caravela]{a caravela — the caravel}

\label{lesson:PT-C328-caravela}



\begin{quote}
Before the new one: say the Portuguese for the crown, then the Portuguese for the navigator.
\end{quote}

\subsection*{You'll want to know: a caravela}
\textbf{a caravela} — \textquotedblleft{}the caravel\textquotedblright{}.

\begin{quote}
A caravela era um barco leve e rápido.
\end{quote}

\begin{quote}
\emph{The caravel was a light, fast ship.}
\end{quote}

The small sailing ship of the age of discovery.

\begin{cousinweb}
\emph{caravela} comes down from Old Portuguese (EtymDB); it is usually traced through Late Latin \emph{carabus} to Greek \emph{kárabos} \textquotedblleft{}a light boat\textquotedblright{}.
\end{cousinweb}

\begin{grammarlens}[title={What you've built}]
One more word for telling, explaining and describing at B1.
\end{grammarlens}

\subsection*{Your turn}
\begin{itemize}
\raggedright
  \item \emph{Say it:} \emph{a caravela}
  \item \emph{Say it:} \emph{a caravela}, once more
  \item \emph{Say it:} a sentence of your own with \emph{a caravela}
  \item \emph{Recall:} say the Portuguese for the crown, then the Portuguese for the navigator, then say \emph{a caravela} again
\end{itemize}

\begin{tcolorbox}[breakable,colback=teal!4,colframe=teal!35!black,title={Before you move on}]
What does a caravela mean? (\textquotedblleft{}The caravel\textquotedblright{}.) Say it once more.
\end{tcolorbox}
\section[naufragar]{naufragar — to be shipwrecked, to sink}

\label{lesson:PT-C328-naufragar}



\begin{quote}
Before the new one: say the Portuguese for the navigator, then the Portuguese for the caravel.
\end{quote}

\subsection*{You'll want to know: naufragar}
\textbf{naufragar} — \textquotedblleft{}to be shipwrecked, to sink\textquotedblright{}.

\begin{quote}
O barco naufragou perto da costa.
\end{quote}

\begin{quote}
\emph{The boat sank near the coast.}
\end{quote}

Also of plans: \emph{o projeto naufragou}.

\begin{cousinweb}
\emph{naufragar} is a learned borrowing from Latin \emph{naufragāre} \textquotedblleft{}to suffer shipwreck\textquotedblright{}, built on \emph{nāvis} \textquotedblleft{}a ship\textquotedblright{} and \emph{frangere} \textquotedblleft{}to break\textquotedblright{} (Lewis \& Short).
\end{cousinweb}

\begin{grammarlens}[title={What you've built}]
One more word for telling, explaining and describing at B1.
\end{grammarlens}

\subsection*{Your turn}
\begin{itemize}
\raggedright
  \item \emph{Say it:} \emph{naufragar}
  \item \emph{Say it:} \emph{naufragar}, once more
  \item \emph{Say it:} a sentence of your own with \emph{naufragar}
  \item \emph{Recall:} say the Portuguese for the navigator, then the Portuguese for the caravel, then say \emph{naufragar} again
\end{itemize}

\begin{tcolorbox}[breakable,colback=teal!4,colframe=teal!35!black,title={Before you move on}]
What does naufragar mean? (\textquotedblleft{}To be shipwrecked, to sink\textquotedblright{}.) Say it once more.
\end{tcolorbox}
\section[o naufrágio]{o naufrágio — the shipwreck}

\label{lesson:PT-C328-naufragio}



\begin{quote}
Before the new one: say the Portuguese for the caravel, then the Portuguese for to be shipwrecked.
\end{quote}

\subsection*{You'll want to know: o naufrágio}
\textbf{o naufrágio} — \textquotedblleft{}the shipwreck\textquotedblright{}.

\begin{quote}
Só dois homens sobreviveram ao naufrágio.
\end{quote}

\begin{quote}
\emph{Only two men survived the shipwreck.}
\end{quote}

\emph{Naufragar} — to be shipwrecked.

\begin{cousinweb}
\emph{naufrágio} is a learned borrowing from Latin \emph{naufragium} \textquotedblleft{}a shipwreck\textquotedblright{} (EtymDB), built on \emph{nāvis} \textquotedblleft{}a ship\textquotedblright{} and \emph{frangere} \textquotedblleft{}to break\textquotedblright{}.
\end{cousinweb}

\begin{grammarlens}[title={What you've built}]
One more word for telling, explaining and describing at B1.
\end{grammarlens}

\subsection*{Your turn}
\begin{itemize}
\raggedright
  \item \emph{Say it:} \emph{o naufrágio}
  \item \emph{Say it:} \emph{o naufrágio}, once more
  \item \emph{Say it:} a sentence of your own with \emph{o naufrágio}
  \item \emph{Recall:} say the Portuguese for the caravel, then the Portuguese for to be shipwrecked, then say \emph{o naufrágio} again
\end{itemize}

\begin{tcolorbox}[breakable,colback=teal!4,colframe=teal!35!black,title={Before you move on}]
What does o naufrágio mean? (\textquotedblleft{}The shipwreck\textquotedblright{}.) Say it once more.
\end{tcolorbox}
\section[o tesouro]{o tesouro — the treasure}

\label{lesson:PT-C328-tesouro}



\begin{quote}
Before the new one: say the Portuguese for to be shipwrecked, then the Portuguese for the shipwreck.
\end{quote}

\subsection*{You'll want to know: o tesouro}
\textbf{o tesouro} — \textquotedblleft{}the treasure\textquotedblright{}.

\begin{quote}
Encontraram um tesouro numa gruta.
\end{quote}

\begin{quote}
\emph{They found a treasure in a cave.}
\end{quote}

Also a term of affection: \emph{o meu tesouro}.

\begin{cousinweb}
\emph{tesouro} is inherited from Latin \emph{thēsaurus} \textquotedblleft{}a treasure\textquotedblright{} (EtymDB), from Greek \emph{thēsaurós} \textquotedblleft{}a storehouse, a treasure\textquotedblright{}.
\end{cousinweb}

\begin{grammarlens}[title={What you've built}]
One more word for telling, explaining and describing at B1.
\end{grammarlens}

\subsection*{Your turn}
\begin{itemize}
\raggedright
  \item \emph{Say it:} \emph{o tesouro}
  \item \emph{Say it:} \emph{o tesouro}, once more
  \item \emph{Say it:} a sentence of your own with \emph{o tesouro}
  \item \emph{Recall:} say the Portuguese for to be shipwrecked, then the Portuguese for the shipwreck, then say \emph{o tesouro} again
\end{itemize}

\begin{tcolorbox}[breakable,colback=teal!4,colframe=teal!35!black,title={Before you move on}]
What does o tesouro mean? (\textquotedblleft{}The treasure\textquotedblright{}.) Say it once more.
\end{tcolorbox}
